\subsection*{Refined products and simultaneous finite-action conjugacy}

\statement{B.46}{Known}{The exact refined-product criterion}
For subgroups $H,J$ of a residually finite group $G$, the
following conditions are equivalent:
\begin{enumerate}
\item $HJ$ is closed in the profinite topology of $G$, and
$\overline H\cap\overline J=\overline{H\cap J}$ in $\widehat G$;
\item for every finite-index subgroup $L\leq G$ containing
$H\cap J$, the subset $(H\cap L)J$ is closed in that topology.
\end{enumerate}
This is \citet[Theorem 1.1]{DE_Minasyan2023}; its proof is
Proposition 2.4 together with Proposition 3.1.
For the actual $P,H_T,H_5$ of B.28, it gives the equivalence
\[
 \left\{\begin{array}{c}
 H_TH_5\text{ is closed in }P,\\
 \mathcal K=1
 \end{array}\right.
 \quad\Longleftrightarrow\quad
 (H_T\cap L)H_5\text{ is closed in }P
 \quad\text{for every }L\leq_f P.
\]
Equivalently, every product $H_0J_0$ is closed for finite-index
$H_0\leq H_T$ and $J_0\leq H_5$.
The quantifier over refinements cannot in general be omitted.
\begin{proof}[Independent proof of the specialization used here]
For this application assume $H\cap J=1$.
All closures below lie in $\widehat G$.
For an open subgroup $L\leq G$,
\[
 \overline{H\cap L}=\overline H\cap\overline L.
\]
Indeed, the right side is open in $\overline H$, and the
dense subgroup $H$ meets it exactly in $H\cap L$.
Also $\overline{XY}=\overline X\,\overline Y$ for two
subgroups $X,Y$: multiplication is continuous, and the
product of the compact closures is closed.

Assume $HJ$ is closed and $\overline H\cap\overline J=1$.
If $g\in G\cap\overline{H\cap L}\,\overline J$, then
$g\in HJ$, so write $g=hj$ with $h\in H,j\in J$.
Write also $g=uv$ with
$u\in\overline{H\cap L},v\in\overline J$.
The equality $h^{-1}u=jv^{-1}$ belongs to
$\overline H\cap\overline J$, so $u=h$ and $h\in L$.
Thus $g\in(H\cap L)J$, proving closure of the refined product.

Conversely assume all those refined products are closed.
The choice $L=G$ gives closure of $HJ$.
Let $a\in\overline H\cap\overline J$, and fix a
finite-index normal subgroup $U\trianglelefteq G$.
Choose $h\in H$ with the same image as $a$ modulo $U$.
Then
\[
 h^{-1}=(h^{-1}a)a^{-1}
 \in\overline{H\cap U}\,\overline J
       =\overline{(H\cap U)J}.
\]
Closure gives $h^{-1}=uj$ with $u\in H\cap U$ and
$j\in J$.  Hence $hu=j^{-1}\in H\cap J=1$, so
$h\in U$ and $a\in\overline U$.  Intersecting over $U$
gives $a=1$.

In B.28 the induced profinite topologies on $H_T$ and $H_5$
are full.  Consequently every finite-index subgroup of either
factor is its intersection with an open subgroup of $P$.
The symmetric version of the argument proves closure of
$H_TJ_0$ as well as $H_0H_5$, and
\[
 (H_0H_5)\cap(H_TJ_0)=H_0(H_T\cap H_5)J_0=H_0J_0.
\]
This proves the last equivalence and matches
\cite[Corollary 3.2]{DE_Minasyan2023}.
\end{proof}

\statement{B.47}{Known}{A closed product with a nontrivial profinite intersection}
Let $V=M_2(\mathbb Z)$ additively, let $S=\mathrm{SL}_2(\mathbb Z)$
act on $V$ by left multiplication, and put
\[
 G=V\rtimes S,
 \qquad H=\{(0,B):B\in S\},
 \qquad J=(I,1)H(I,1)^{-1}.
\]
Then $G$ is residually finite, $H\cap J=1$, and $HJ$ is
closed in the full profinite topology of $G$.
Both factors have the full induced profinite topology; every
finite-index subgroup of either factor is separable in $G$.
Nevertheless
\[
 \overline H\cap\overline J
 \cong\ker\bigl(\widehat{\mathrm{SL}_2(\mathbb Z)}
          \longrightarrow\mathrm{SL}_2(\widehat{\mathbb Z})\bigr)
 \ne1.
\]
In fact this intersection has a quotient isomorphic to
$\mathrm{PSL}_2(\mathbb F_8)$.
This is the example of \citet[Example 4.1]{DE_Minasyan2023},
with the failure witnessed here by an explicit finite quotient.
It is a counterexample to the general implication from one
closed product to equality of intersections.  It is not a
counterexample to the icosahedral assertion.
\begin{proof}
Our multiplication convention is
\[
 (A,B)(A',B')=(A+BA',BB').
\]
Thus $J=\{(I-B,B):B\in S\}$ and
\[
 HJ=\{(A,B)\in V\rtimes S:\det(A+B)=1\}.
\]
The two factors intersect only if $I-B=0$, hence only at the
identity.  If $(A,B)$ is outside the displayed product, choose
an integer $m\geq2$ which does not divide
$\det(A+B)-1$.  Reduction to the finite group
\[
 M_2(\mathbb Z/m)\rtimes\mathrm{SL}_2(\mathbb Z/m)
\]
separates it from $HJ$.  These congruence quotients also
separate every nonidentity element of $G$, proving residual
finiteness.

The retraction $G\to S$, $(A,B)\mapsto B$, gives the full
induced topology on $H$.  The subgroup $H$ itself is closed:
a nonzero translation coordinate remains nonzero modulo some
$m$.  If $H_0\leq_f H$, intersect $H$ with the inverse image
of $H_0$ under this retraction.  This writes $H_0$ as the
intersection of a closed subgroup with an open subgroup of
$G$, proving its separability.  Conjugation gives the same
facts for $J$.

There is a topological-group isomorphism
\[
 \widehat G\cong
 M_2(\widehat{\mathbb Z})\rtimes\widehat S,
\]
where the action factors through
$\widehat S\to\mathrm{SL}_2(\widehat{\mathbb Z})$.
Here is the full-topology check.  In any finite quotient of
$G$, the image of $V$ has finite exponent, so some $mV$ is
killed.  Intersect the kernel of the restricted map on $S$
with its kernel on $V/mV$.  This is a finite-index normal
subgroup $W$ of $S$, and the original quotient factors through
$(V/mV)\rtimes(S/W)$.  Such quotients are cofinal, proving
the asserted completion.  Reduction of integral elementary
matrices, followed by the Chinese remainder theorem, shows
that the congruence map onto
$\mathrm{SL}_2(\widehat{\mathbb Z})$ is surjective.

Inside this semidirect product an element $(0,b)$ belongs to
both closures exactly when $bI=I$.  Since the action is left
matrix multiplication, this means that the congruence image
of $b$ is the identity matrix.  This proves the kernel formula,
without assuming the kernel is trivial.

We now verify its nontriviality by a finite quotient.  Write
$\mathbb F_8=\mathbb F_2[u]/(u^3+u+1)$ and take
\[
 a=\begin{pmatrix}1&1\\0&1\end{pmatrix},
 \qquad
 b=\begin{pmatrix}0&u^{-1}\\u&1\end{pmatrix}.
\]
Exact enumeration in
\texttt{check\_separability\_example.py} proves that $a,b,ab$
have orders $2,3,7$, that $a,b$ generate all $504$ determinant-one
matrices, and that the group is simple.  For the last assertion,
the program enumerates its nine conjugacy classes, of sizes
\[
 1,56,56,56,56,63,72,72,72,
\]
and verifies that the normal closure of every nonidentity
class is the whole group.  Characteristic two makes its center
trivial.  Thus the presentation
$\mathrm{PSL}_2(\mathbb Z)=C_2*C_3$ gives a surjection
\[
 S\longrightarrow Q=\mathrm{PSL}_2(\mathbb F_8),
\]
which kills the central involution of $S$.

This $Q$ is not a quotient of
$\mathrm{SL}_2(\widehat{\mathbb Z})$.
Indeed any continuous finite quotient factors through
$\mathrm{SL}_2(\mathbb Z/m)$.  By Chinese remaindering and
simplicity, a nonabelian simple quotient comes from one
prime-power factor.  The kernel of reduction from
$\mathrm{SL}_2(\mathbb Z/p^k)$ to
$\mathrm{SL}_2(\mathbb F_p)$ is a $p$-group and must have
trivial image in a nonabelian simple group.  Removing the
center shows that the simple quotient is
$\mathrm{PSL}_2(\mathbb F_p)$ for a prime $p\geq5$;
the factors for $p=2,3$ are solvable.
The orders at $p=5,7,11$ are respectively $60,168,660$,
and the formula $p(p^2-1)/2$ is strictly increasing.
None has order $504$.

Consequently the surjection $\widehat S\to Q$ does not
kill the congruence kernel.  Its image on that kernel is
normal and nontrivial in $Q$, hence all of $Q$.
This proves the final assertion.  The exact program and its
standard output are included; no floating-point or
unverified finite-group identification is used.
\end{proof}

The failure persists with torsion-free free factors.
Put $\Lambda=\ker(S\to\mathrm{SL}_2(\mathbb F_3))$ and
$G_3=V\rtimes\Lambda$, and intersect $H,J$ with $G_3$.
Both resulting factors are free of rank three.
Indeed $-I\notin\Lambda$, so $\Lambda$ injects into
$\mathrm{PSL}_2(\mathbb Z)=C_2*C_3$.
Its image is the kernel of reduction to
$\mathrm{PSL}_2(\mathbb F_3)$ and has index twelve.
The usual generators represented by
$\left(\begin{smallmatrix}0&-1\\1&0\end{smallmatrix}\right)$
and $\left(\begin{smallmatrix}0&-1\\1&1\end{smallmatrix}\right)$
remain of orders two and three in that quotient.
Reduction is therefore injective on both finite factors.
Every finite-order element of their free product is conjugate
into a factor, so the kernel is torsion-free and free.
Its free rank is $1-12(1/2+1/3-1)=3$.
In $G_3$ the product is the closed set
\[
 \{(A,B):\det(A+B)=1,\quad A+B\equiv I\pmod3\}.
\]
The congruence kernel above lies in $\widehat\Lambda$,
since it is killed by reduction modulo three.
The completion of the finite-index subgroup $G_3$ is its
open closure in $\widehat G$.
Intersecting the two previous factor closures with this
open subgroup therefore leaves their intersection unchanged.
Thus the two closures still have the same nontrivial
intersection.  Retraction and conjugation again give
full induced topologies and separability of every
finite-index subgroup of the factors.
This strengthens the general counterexample without
making any claim about the actual symmetric subgroups.

\statement{B.48}{Known}{The hypotheses of the recent separability criteria}
Put $\Pi=\pi_1(S^2\setminus V_{12})\cong F_{11}$ and use
the faithful Dehn--Nielsen map $P\hookrightarrow\operatorname{Out}(\Pi)$.
For $H\leq P$, write
\[
 \widetilde H=\{\alpha\in\operatorname{Aut}(\Pi):
                        [\alpha]\in H\}.
\]
If $\widetilde H$ is conjugacy separable, then $H$ is separable
in $\operatorname{Out}(\Pi)$, and therefore in $P$.
This is the exact specialization of
\cite[Corollary D]{DE_FournierFacio2025}.
For $H=H_T,H_5$ the preimages are, after choosing lifts of
free bases, $F_{11}\rtimes F_{13}$ and $F_{11}\rtimes F_6$.
The cited criterion concerns individual subgroups and does
not assert closure of their products.
\begin{proof}[Application and hypothesis audit]
The group $\Pi$ is torsion-free hyperbolic.
The cited corollary applies to every subgroup of
$\operatorname{Out}(\Pi)$.
Its hypotheses do not include convex cocompactness.
In the source, Lemma 2.3 produces an inner automorphism whose
centralizer in $\operatorname{Aut}(\Pi)$ is cyclic.
Corollary 2.2 uses that element to distinguish an outer
automorphism outside $H$ from the $\widetilde H$ orbit.
Proposition 2.1 then uses conjugacy separability of
$\widetilde H$ to obtain the required finite quotient.
Thus the extension-group hypothesis remains essential at
the precise finite-separation step.  The exact sequence
\[
 1\longrightarrow\Pi\longrightarrow\widetilde H
       \longrightarrow H\longrightarrow1
\]
splits when $H$ is free, by choosing lifts of its free
generators.  This proves the two semidirect-product
descriptions; their conjugacy separability is not inferred
from the freeness of the two factors.

For the actual centralizers, individual separability already
has an elementary proof.  If $g\in P\setminus C_P(T)$,
some $t\in T$ has $[g,t]\ne1$ in the residually finite
group $\Gamma_{0,[12]}$.  A finite quotient detecting this
commutator separates $g$ from $H_T$.  The same argument
works for $H_5$.  The full induced topologies of B.28
then make every finite-index subgroup of each separable:
given $H_0\leq_f H$, choose $U\trianglelefteq_f P$
with $H\cap U\subseteq H_0$, and use
$H_0=H\cap H_0U$.  B.47 shows why these individual
facts do not settle B.46.

The earlier result is
\citet[Theorem 1.1]{DE_HagenSisto2023}.
It concerns a subgroup of a closed-surface mapping class group.
The subgroup must be torsion-free, malnormal and convex cocompact,
and its Birman preimage must be conjugacy separable.
Those geometric
hypotheses are not supplied by B.28.  In fact the symmetric
subgroups here contain nontrivial multitwists, and so do
their finite-index subgroups.  To see this directly for
$T$, take on its four-point quotient sphere a curve enclosing
the free puncture and the cone point of order two.  Its
preimage has six components, each bounding two of the
twelve punctures; a suitable power of its twist lifts to
the product of twists about these six essential curves.
For $D$, a quotient curve enclosing one free puncture and
one fixed pole lifts to an essential curve separating six
marks from the other six.  A power of the quotient twist
lifts to its nontrivial twist.  These lifts are pure and
commute with the corresponding deck group.  A nontrivial
multitwist has bounded orbit in the curve graph, whereas
the orbit map of a convex-cocompact subgroup is a
quasi-isometric embedding; see the characterization in
\cite[Section 2.1]{DE_HagenSisto2023}.  Passing to finite
index preserves a nonzero power of each such element.

The product theorem is
\cite[Corollary 1.4]{DE_MinasyanMineh2025}.
Its ambient group must be finitely generated and relatively
hyperbolic.  Every finitely generated relatively quasiconvex
subgroup must be separable, the peripheral subgroups must
be double-coset separable, and both factors must be
relatively quasiconvex.
Its proof in Section 15 forms a relatively quasiconvex
join of finite-index subgroups and separates an unwanted
element from finitely many translates of that join.
No suitable relatively hyperbolic structure has been
proved here for $P$.
Nor have we found a common ambient subgroup of $P$
with those properties and the required induced profinite
topology.  Thus that theorem has not been applied to
the products in B.46.
The central-extension result of
\citet[Theorem 1.3]{DE_Mineh2026} requires a
double-coset-separable quotient, a finitely generated
central kernel, and subgroup separability of the total
group.  We have not supplied such a comparison for an
ambient group containing both symmetric subgroups.
A braid central extension by itself does not verify
the condition on its quotient.
\end{proof}

\statement{B.49}{Proved}{A finite-action reformulation of every refined product}
Retain the fixed embedding $\iota:A_5\hookrightarrow
\Gamma=\Gamma_{0,[12]}$, with the generating subgroups
$T=A_4$ and $D=C_5$, and identify them with their images.
Let $L\leq_f P$ and $g\in P$.
Then $g$ belongs to the closure of $(H_T\cap L)H_5$ in
$\widehat P$ if and only if there are a monomorphism
$\iota_g:A_5\hookrightarrow\Gamma$ and
$u\in\overline L\leq\widehat P$ satisfying
\[
 \iota_g(t)=t\quad(t\in T),\qquad
 \iota_g(d)=g d g^{-1}\quad(d\in D),\qquad
 \iota_g(a)=u a u^{-1}\quad(a\in A_5).
\]
The monomorphism, when it exists, is unique from its first
two prescriptions.  For such a monomorphism,
\[
 g\in(H_T\cap L)H_5
 \quad\Longleftrightarrow\quad
 \iota_g=h\iota h^{-1}\text{ for some }h\in L.
\]
Consequently the refined closure condition in B.46 is
exactly the following assertion, for all $L$ and $g$:
every displayed homomorphism conjugacy realized in
$\overline L$ is realized in $L$.
This is simultaneous conjugacy of homomorphisms, with
the restrictions on $T$ and $D$ retained.
\begin{proof}
The topology induced on $P$ from $\Gamma$ is full because
$P$ has finite index, so we regard $\widehat P$ as an open
subgroup of $\widehat\Gamma$.
Suppose
\[
 g=uv,
 \quad u\in\overline{H_T\cap L},
 \quad v\in\overline{H_5}.
\]
Conjugation by $u$ fixes $T$ pointwise, and on $D$ it is
conjugation by $g$, because $v$ centralizes $D$.
Thus $uA_5u^{-1}$ is generated by the two actual finite
subsets $T$ and $gDg^{-1}$ of $\Gamma$.
Every one of its elements is therefore an element of
$\Gamma$, not just of its completion.  Conjugation by $u$
gives the claimed monomorphism $\iota_g$ into that
discrete subgroup.  All the $A_5$ relations hold there
because $\Gamma\hookrightarrow\widehat\Gamma$ is
injective.  This is the step ensuring that the proposed
representation has finite image; it is not an assumption
about arbitrary subgroups generated by two finite groups.

Conversely let $\iota_g$ and $u$ satisfy the prescriptions.
Since $u$ centralizes $T$, B.28 gives
$u\in\overline{H_T}$.  Because $L$ is open,
\[
 u\in\overline{H_T}\cap\overline L
       =\overline{H_T\cap L}.
\]
The equality of conjugations on $D$ says that
$u^{-1}g\in C_{\widehat P}(D)=\overline{H_5}$,
again by B.28.  This gives the required product
factorization in the completion.

The same argument in the discrete group proves the second
equivalence.  A conjugator $h\in L$ fixes $T$ pointwise,
so it is in $H_T\cap L$, and $h^{-1}g$ centralizes $D$,
so it is in $H_5$.  The reverse direction follows by
conjugating with the first product factor.
Finally, the two prescriptions determine a homomorphism
uniquely because $\langle T,D\rangle=A_5$.
This proves all the equivalences.
\end{proof}

For $L=P$, B.52 below proves the required conjugacy
and closure of the actual unrefined product.
The assertion for every finite-index $L$ remains unresolved.
A theorem for single torsion elements or solvable finite
subgroups does not give that marked homomorphism comparison.
Nor has a profinite complex with the required unique
$A_5$ fixed point been constructed.
The quantified refinement remains the precise unresolved
form of route B.4(c), and $\mathcal K=1$ remains Open.
